\section{Runtime Control under Qualified Fixed Modes}
\label{sec:runtime}
The separate M4 runtime study tests whether continuation control improves a qualified speculative path. It uses its own frozen policy, direct-ID gate and cost boundary; historical M4 timing is background rather than its qualification evidence.

\subsection{A prospective cost guard on M4}
Guard tests a frozen rule against target-only and fixed speculation, using Qwen3-8B/0.6B Q8\_0 on llama.cpp Metal, fixed $k=\RTK$, greedy generation and a \RTCap-token cap. It starts on Always-Spec's path and may fall back permanently to target-only. At eligible completed rounds it divides all charged generation cost by committed progress, including the initial seed and correction/bonus; rejected proposals do not count. It continues when cost per token is at most $(1-m)R_{\mathrm{ref}}$, otherwise falling back after reconciling invalid suffix state. Zero progress also triggers fallback.

Calibration freezes the target reference and selects margin \RTMargin\ with decisions starting after \RTRounds\ rounds. The controller receives no task labels, future length or request-specific oracle. The post-prefill timer charges control and transition costs; Appendix T retains the full rule, grid, tie-breaks and timing boundary.

\subsection{Held-out evaluation}
Historical outcomes informed the design; prospective means policy freeze before reserved outcomes. Exposure exclusion and deterministic selection give disjoint calibration, qualification and held-out splits of \RTCalibration, \RTQualification\ and \RTPrompts\ prompts, balanced across HumanEval and Spec-Bench. Policy and reference precede qualification/formal measurement; the appendix retains the audited freeze-hash metadata defect. Always-Spec and Guard each qualify \RTQualPairs/\RTQualPairs\ direct-ID pairs.

The held-out study runs \RTReps\ repetitions of three conditions on \RTPrompts\ prompts (\RTObservations\ observations), randomizing condition order within prompt/repetition. Both speculative paths preserve \RTExact/\RTExact\ complete direct-ID sequences. Prompt condition medians precede latency ratios and the GM; \RTBootstrap\ prompt-bootstrap draws use seed \RTSeed. The supplement retains exact estimator, exposure and metadata-recovery details.

\begin{table}[t]
\caption{M4 per-prompt median post-prefill latencies. Each speculative condition is \RTExact/\RTExact\ direct-token exact. The offline oracle is unavailable online. The separated Guard/Always-Spec contrast is reviewer-requested post-study paired secondary analysis; values above one mean Guard is slower.}
\label{tab:runtime}
\centering
\begin{tabular}{lcc}
\toprule
Latency ratio & GM & 95\% prompt CI \\
\midrule
\multicolumn{3}{l}{\emph{Speedup versus target-only (higher is faster)}}\\
Target / Guard & \RTGuard & [\RTGuardLo, \RTGuardHi] \\
Target / Always-Spec & \RTAlways & [\RTAlwaysLo, \RTAlwaysHi] \\
\midrule
\multicolumn{3}{l}{\emph{Descriptive hindsight latency ratio}}\\
Guard / offline oracle & \RTOracle & [\RTOracleLo, \RTOracleHi] \\
\midrule
\multicolumn{3}{l}{\emph{Post-study paired secondary latency ratio}}\\
Guard / Always-Spec & \RTSecondary & [\RTSecondaryLo, \RTSecondaryHi] \\
\bottomrule
\end{tabular}
\end{table}

Guard is faster than target-only, but fixed speculation has the higher speedup (Table~\ref{tab:runtime}). A post-study paired comparison gives Guard/Always-Spec latency \RTSecondary\ [\RTSecondaryLo, \RTSecondaryHi], entirely above parity. Guard is faster on \RTGuardFaster\ prompt medians, Always-Spec on \RTAlwaysFaster, with \RTTies\ ties. Fallback occurs in \RTFallback\% of \RTExact\ executions, measuring activity rather than beneficial decisions. Guard/oracle is \RTOracle$\times$; this hindsight oracle chooses the lower target-only or Always-Spec median, not an achievable runtime policy.

Reconstructing all original medians finds target-only faster than Always on only \DerTargetFaster/\RTPrompts\ prompts, both Spec-Bench. Always/oracle is \DerAlwaysOracle$\times$. Choosing the better of these two measured fixed modes after seeing outcomes therefore offers about \DiagMOpportunity\% GM latency reduction from Always, computed from the full-precision ratio. Guard falls back in all six executions of those two prompts but still loses to target-only on their medians. The family-stratified sensitivity analysis also leaves Guard slower (Appendix P). These post-study comparisons estimate neither transition cost nor opportunity for arbitrary within-request switching.

Guard preserves IDs but supplies no incremental benefit over Always-Spec. Fallback activity does not establish beneficial decisions or risk reduction. No runtime energy claim is made. Policy transfer and stable small-stratum effects remain unestablished; the archived policy-freeze defect and original stopped pipeline markers are documented separately.
